% Literature-aligned research-paper skeleton for the DNS timing-sequence thesis.
%
% Design principle:
% Keep the paper close to the compact structure repeatedly used in DNS-tunneling
% and covert-timing-channel research: Introduction -> Related Work -> System /
% Experimental Design -> Method -> Evaluation/Results -> Discussion -> Conclusion.
%
% All prose is intentionally blank. Comments are writing guidance only and do
% not appear in the compiled paper. Figure/table blocks remain commented until
% the corresponding artifact is ready.

\documentclass{vgtc}                          % final (conference style)
%\documentclass[review]{vgtc}                 % review
%\documentclass[widereview]{vgtc}             % wide-spaced review
%\documentclass[preprint]{vgtc}               % preprint
%\documentclass[electronic]{vgtc}             % electronic version

\ifpdf
  \pdfoutput=1\relax
  \pdfcompresslevel=9
  \pdfoptionpdfminorversion=7
  \ExecuteOptions{pdftex}
  \usepackage{graphicx}
  \DeclareGraphicsExtensions{.pdf,.png,.jpg,.jpeg}
\else
  \ExecuteOptions{dvips}
  \usepackage{graphicx}
  \DeclareGraphicsExtensions{.eps}
\fi

\graphicspath{{figures/}{pictures/}{images/}{./}}

\usepackage{microtype}
\PassOptionsToPackage{warn}{textcomp}
\usepackage{textcomp}
\usepackage{mathptmx}
\usepackage{times}
\renewcommand*\ttdefault{txtt}
\usepackage{cite}
\usepackage{tabu}
\usepackage{booktabs}
\usepackage{enumitem}
\usepackage{setspace}

\onlineid{0}
\vgtccategory{Research}
\vgtcinsertpkg

\title{Detecting DNS Timing-Sequence Data Exfiltration with Benign-Calibrated Temporal Anomaly Detection}

\author{Sotiris Tofinis\\ %
        \scriptsize Open University of Cyprus
}

% Write the abstract last:
% problem -> method -> real benign data + controlled attacks -> main results ->
% measured limitation / contribution.
\abstract{}

\begin{document}

\firstsection{Introduction}
\maketitle

% ======================================================================
% 1. INTRODUCTION
% ======================================================================
%
% Keep this as one continuous section, as is common in the papers reviewed.
% Suggested paragraph flow:
%   1. DNS is necessary and therefore attractive for covert exfiltration.
%   2. Timing-sequence exfiltration carries information in WHEN queries occur,
%      so content/signature inspection is not sufficient for this threat.
%   3. Research gap: detecting weak timing anomalies against real human endpoint
%      DNS behavior, using benign-only and content-free modeling.
%   4. State the aim / research questions in prose.
%   5. End with 3--4 explicit contributions.


% ======================================================================
% 2. BACKGROUND AND RELATED WORK
% ======================================================================

\section{Background and Related Work}

\subsection{DNS Exfiltration and Covert Timing Channels}
% Explain DNS exfiltration briefly, then narrow immediately to timing channels.
% Cover inter-arrival times, low-and-slow behavior, jitter/evasion, and the
% difference between content-based DNS tunneling and timing-sequence exfiltration.

% FIGURE 1 -- DNS timing-channel concept.
% One simple timeline: benign DNS activity vs. timing-encoded DNS activity.
% The visual should make the core thesis understandable without equations.
%\begin{figure}[t]
%  \centering
%  \includegraphics[width=\columnwidth]{fig01_dns_timing_channel_concept}
%  \caption{Conceptual example of information encoded in the timing sequence of DNS requests.}
%  \label{fig:timing-concept}
%\end{figure}

\subsection{Existing Detection Approaches and Research Gap}
% Organize related work by approach rather than one subsection per paper:
% statistical timing detectors -> ML timing detectors -> DNS tunneling/exfiltration
% detectors -> benign/unsupervised behavioral detectors.
% End with the exact gap addressed by this thesis.

% TABLE 1 -- Compact related-work positioning table.
% Suggested columns:
% Work | DNS | Timing-aware | Benign/unsupervised | Content-free | Real benign traffic.


% ======================================================================
% 3. SYSTEM AND EXPERIMENTAL DESIGN
% ======================================================================

\section{System and Experimental Design}

\subsection{Data, Telemetry, and Testbed}
% In one subsection cover:
%   - real benign endpoint DNS capture and duration;
%   - endpoint/unicast-DNS filtering and Zeek telemetry;
%   - attack lab topology and capture point;
%   - parent-domain grouping, 15-minute inactivity/session policy;
%   - threat/scope statement: endpoint timing detector, not payload inspection.

% FIGURE 2 -- Experimental design / data pipeline.
% Combine topology and processing into one useful figure:
% real benign endpoint + controlled attack generator -> Zeek DNS events ->
% filtering/grouping -> causal feature windows -> detector evaluation.
%\begin{figure}[t]
%  \centering
%  \includegraphics[width=\columnwidth]{fig02_experimental_pipeline}
%  \caption{Experimental setup and DNS telemetry processing pipeline.}
%  \label{fig:experimental-pipeline}
%\end{figure}

\subsection{Timing Attack Profiles}
% Describe the attack family here, not in many separate subsections:
% no_jitter; J1--J3; S1 pair-order; S2 context-coupled low-rate; and the S2
% 0.50, 0.25, and 0.125 boundary variants.
% Explain what each family was designed to test and verify payload integrity.

% TABLE 2 -- Dataset and attack-profile specification.
% Suggested columns:
% Profile | Timing mechanism | Queries | Duration | Purpose | Payload verified.

% FIGURE 3 -- Attack timing profiles.
% A compact multi-panel/timeline view of no_jitter, J1--J3, S1 and the S2 family.
% Emphasize how S2 becomes progressively slower / more blended with benign timing.


% ======================================================================
% 4. DETECTION METHOD
% ======================================================================

\section{Detection Method}

\subsection{Causal Timing Representation}
% Explain only what the reader needs to reproduce the inputs:
% trailing same-parent-domain windows; 64-event current window; adjacent prior
% window for local shifts; no future leakage; undefined first IATs remain null.
% Then define the final 13 content-free timing features as a table rather than
% giving every feature its own prose subsection.

% TABLE 3 -- Final 13 timing features.
% Suggested columns: Feature | Definition | Behavioral interpretation.

\subsection{Locked Direct Detector}
% Robust benign scaling -> regularized covariance -> squared Mahalanobis score ->
% benign-calibrated q99.9 locked threshold. State explicitly that attacks do not
% train the detector and that the direct threshold is not lowered for S2.

\subsection{Sequential Evidence Extension}
% Explain why this second path was added: 0.25/0.50 S2 produced persistent weak
% anomaly evidence but no individual direct-threshold crossing.
%
% Present the three equations together:
%   empirical benign upper-tail probability p_t
%   e_t = max(0, ln(p_0 / p_t))
%   A_t = exp(-Delta t / tau) A_prev + e_t
% Then give the benign-calibrated sequential threshold and combined OR rule.
% Mention 15-minute incident grouping here.

% FIGURE 4 -- Final two-path detector architecture (essential).
% Timing features -> locked Mahalanobis score -> direct threshold
%                          |
%                          -> benign tail probability -> evidence accumulator
%                             -> sequential threshold
% direct OR sequential -> alert -> incident grouping.
%\begin{figure}[t]
%  \centering
%  \includegraphics[width=\columnwidth]{fig04_detector_architecture}
%  \caption{Locked direct detector with the benign-calibrated sequential evidence extension.}
%  \label{fig:detector-architecture}
%\end{figure}


% ======================================================================
% 5. EVALUATION AND RESULTS
% ======================================================================

\section{Evaluation and Results}

\subsection{Evaluation Protocol and Benign Behavior}
% Define training/calibration/validation roles, eligible windows, locked parameters,
% detection criteria, detection delay, and benign incidents/day.
% State clearly which attack profiles were used during detector development.

\subsection{Direct Detector and the S2 Sensitivity Boundary}
% First present the locked direct results for all attacks.
% Then focus on the central boundary experiment:
% 0.125 crosses the direct threshold; 0.25, 0.50 and original S2 do not.
% Use normalized score/threshold so the boundary is visually obvious.

% FIGURE 5 -- Locked direct results + S2 boundary (essential).
% Prefer a two-panel figure:
%   A. all attack profiles vs. direct threshold;
%   B. only the S2 family, ordered by rate/scale, with y=1 boundary.

\subsection{Expanded Detector Results}
% Present the direct-vs-expanded comparison and the key finding:
% sequential evidence recovers 0.25 and 0.50 (and detects 0.125 earlier), while
% original S2 remains below threshold and defines the measured low-rate limit.
% Report detection delay and benign incident rate here as part of the tradeoff.

% FIGURE 6 -- S2 sequential-evidence trajectories (essential).
% Four panels: 0.125, 0.25, 0.50, original S2.
% Plot direct score/threshold and accumulator/threshold over attack progress.
% The original-S2 panel should clearly approach but not cross the boundary.

% TABLE 4 -- Main results table; keep all headline numbers in one place.
% Suggested columns:
% Profile | Direct detected | Expanded detected | First path | Delay |
% Max direct/threshold | Max sequential/threshold.
% Add benign validation incident rate immediately below or in a separate small row.


% ======================================================================
% 6. DISCUSSION AND LIMITATIONS
% ======================================================================

\section{Discussion and Limitations}
% Keep this as one coherent section unless it becomes too long.
%
% Suggested flow:
%   1. Answer the research questions directly from the results.
%   2. Explain why direct scoring detects strong/local timing anomalies while the
%      sequential path detects persistent weak evidence in slower S2 variants.
%   3. Interpret the speed--stealth tradeoff and original S2 as the measured
%      sensitivity/evasion boundary of the current detector.
%   4. Give operational meaning: real human endpoint, content-free features,
%      benign-only calibration, detection delay, and incident rate.
%   5. Limitations: one benign endpoint; controlled attack lab; attack-informed
%      development of the sequential extension; original S2 miss; concept drift;
%      no claim of malicious intent/payload decoding; limited mechanism-attribution
%      controls for S1; need independent endpoints/attacks for broad generalization.


% ======================================================================
% 7. CONCLUSION
% ======================================================================

\section{Conclusion and Future Work}
% One concise section:
% restate the problem -> final detector -> strongest empirical result -> honest
% sensitivity limit -> next steps (additional endpoints, unseen captures, repeated
% seeds/start times, matched timing-null controls, concept drift/live deployment).


% ======================================================================
% REFERENCES
% ======================================================================

%\bibliographystyle{abbrv}
\bibliographystyle{abbrv-doi}
%\bibliographystyle{abbrv-doi-narrow}
%\bibliographystyle{abbrv-doi-hyperref}
%\bibliographystyle{abbrv-doi-hyperref-narrow}

\bibliography{template}

\end{document}
